%% sec-6-mobile-deployment.tex -- Mobile Deployment

\section{Mobile Deployment}
\label{sec:mobile}

\subsection{Current State of the Art}
The Snapdragon~8~Gen~4 + Adreno~830 is among the strongest mobile
SoCs in H1~2026, delivering around 2~TFLOPs of FP32 throughput and
a mature Vulkan~1.3 driver. Representative deployed 4DGS works
include:
\begin{itemize}
  \item Mobile-GS~\cite{du2026mobilegs} and Flux-GS~\cite{du2026fluxgs}:
        real-time desktop 3DGS rasterisation at 127--147~FPS on
        Snapdragon~8~Gen~3; the 4D extension is still exploratory but
        the rasteriser kernel and tile-sort logic transfer directly.
  \item Lumina~\cite{feng2025lumina} (ISCA~2025, SJTU + Rochester):
        a mobile neural-rendering benchmark and pipeline that
        exploits computational redundancy for $4.5\times$ speed-up
        and $5.3\times$ energy efficiency.
  \item GS-NFS~\cite{ghosh2026gsnfs} (NVIDIA Research): bandwidth-
        adaptive streaming of dynamic Gaussian splats and point
        clouds at 25~FPS decode on a Jetson~Orin edge GPU, the
        closest hardware analogue to a Snapdragon-class target.
  \item AirGS~\cite{wang2025airgs} (ACM SIGGRAPH~2025): real-time
        4D Gaussian streaming for free-viewpoint video, with
        GPU-friendly inter-frame delta encoding.
  \item PD-4DGS~\cite{li2026pd4dgs}: progressive decomposition of
        4DGS for bandwidth-adaptive streaming, fits a typical
        5G uplink budget by transmitting only the dynamic residual.
  \item StreamSTGS~\cite{ke2025streamstgs} (CVPR~2025): streaming
        spatial + temporal Gaussian grids for real-time
        free-viewpoint video, with a tetrahedral spatial index.
  \item EvoGS~\cite{shi2026evogs}: continuous-layered Gaussian
        splatting organised by an evolution tree, enabling
        scalable progressive 3D streaming on bandwidth-limited
        clients.
  \item ZipSplat~\cite{veicht2026zipsplat}: prunes Gaussians by
        importance scoring, achieving comparable PSNR with
        $\sim$2--3$\times$ fewer splats---directly relevant for
        mobile fillrate budgets.
  \item CodecSplat~\cite{yu2026codecsplat}: ultra-compact latent
        coding on top of feed-forward 3DGS, the lowest model
        size in the surveyed literature (sub-MB indoor scenes).
  \item P-4DGS~\cite{wang2025p4dgs} (CVPR~2025): predictive 4DGS
        with $90\times$ compression via motion-prediction
        residuals---the canonical "predict-then-residual" pattern.
  \item GIFStream~\cite{li2025gifstream} (CVPR~2025): 4DGS-based
        immersive video with a per-frame feature stream, designed
        for VR/AR walkthroughs.
  \item 4DGCPro~\cite{zheng20254dgcpro}: hierarchical 4D Gaussian
        compression for progressive volumetric-video streaming,
        with octree-coherent layer ordering.
\end{itemize}

\subsection{Core Bottlenecks}
\begin{enumerate}
  \item \textbf{Memory bandwidth.} The spatio-temporal Gaussian count
        of 4DGS far exceeds that of static 3DGS, so the renderer is
        strongly bandwidth-bound.
  \item \textbf{Vulkan~1.3 driver maturity.} Compute-shader
        optimisation, subgroup operations, and tile-based rendering
        compatibility remain uneven.
  \item \textbf{Power wall.} Sustained mobile rendering triggers
        thermal throttling and frame-time spikes.
  \item \textbf{Storage of the dynamic part.} Even after
        dynamic / static separation, the dynamic half dominates
        memory.
\end{enumerate}

\subsection{Project Direction}
The companion \texttt{README.md} of this project fixes the target as
\emph{real-time dynamic 4DGS on Snapdragon~8~Gen~4 / Adreno~830 +
Vulkan~1.3}; see \texttt{01-} / \texttt{02-} / \texttt{03-} for the
engineering plan.

\bigskip
\noindent\textbf{Index of paper notes backing this section}:
\texttt{INDEX.md} \S D / E.
